We begin by describing our construction. For any $n$, $\alpha < n$ and $d> \sqrt{n/\alpha}$, we construct a graph $G = (V, E)$ over $n$ vertices, average degree $d$ and arboricity $\alpha$, such that the following holds.
For any $v\in V$, if we take $q = o()$ length-$2$ random walks from $v$, then with high probability we will not see a $C_4$.
Moreover, $G$ is far from being $C_4$-free. 


\paragraph{The construction}
We shall use the following bipartite graph in one of our lower bounds constructions. 
In this graph, we view one side as corresponding to “lines” and the other to “points”.
So if we view $(a,b)$ as the parameters of a line, then its neighbors are all points $(x,y)$ (in $Z_p \times Z_p$) for which $x = b y  - a$. 

\begin{definition}
[The bipartite graph $G^p_F$]
Let $p$ be a positive integer.
The vertex set of $G^p_F$ is $V_1\cup V_2$ where both in $V_1$ and $V_2$ we have vertices labeled by $Z_p \times Z_p$. There is an edge between $(a,b)\in V_1$ and $(x,y)\in V_2$ iff $x+ a = b\cdot y$ (over $Z_p$). 
\end{definition}

We begin with three disjoint subsets of vertices $V_1$, $V_2$, and $V_3$, each of size $n/3$ (for the sake of brevity we assume that $n$ is divisible by $3$).
We define the following tripartite graph. 
The subgraph induced on $A\cup B$ is $G_F$ (see Definition~\ref{}) where $A$ is one part and $B$ is the other part. Likewise for the subgraph induced on $B\cup C$ w.r.t. $B$ and $C$.
By the construction of $G_F$, at this point for any pair of distinct vertices $u\in B$, $v\in B$, there is a single vertex in $A$ that is adjacent to both vertices and likewise for $C$.
Therefore there is a unique $C_4$ which includes both $u$ and $v$.

We next remove some edges from the graph. For each vertex $v\in B$, with probability $1- p_{\rm ver}$ where $p_{\rm ver} = \alpha/d$, we remove all the edges adjacent to $v$. We then independently remove every edge in the resulting graph with probability $1- p_{\rm e}$ where $p_{\rm e} = d/\sqrt{n}$. This completes the description of the construction. 

\begin{claim}
The graph is far from being $C_4$-free.
\end{claim}
\begin{proof}
    The expected number of $C_4$s is:
    \begin{equation}
        {|V_2| \choose 2} \cdot \left(\frac{\alpha}{d}\right)^2 \left(\frac{d}{\sqrt{n}}\right)^4 = \Theta(\alpha^2 d^2)
    \end{equation}
    Let $(x, y) \in (V_1\times V_2) \cup (V_3\times V_2)$. The expected number of paths of length $3$ from $x$ to $y$ is bounded by 
    \begin{equation}
        \sqrt{n} \cdot p^2_{\rm ver} \cdot p^3_{\rm e} = \theta(\frac{\alpha^2 d}{n}).
    \end{equation} 
\end{proof}
In fact, both hold not only in expectation, but with high probability. This holds by Chebyshev's inequality. To see this, for each pair in $(a, b)\in {V_2 \choose 2}$ consider a random variable $Z_{a,b}$ which is an indicator for the event that there is a $C_4$ which includes both $a$ and $b$ in $G$.
For $a \neq a'$ and $b \neq b'$ it holds that $Z_{a,b}$ and $Z_{a',b'}$ are independent.
Therefore we would like to bound the contribution of $Z_{a,b}$ and $Z_{a,b'}$ to the variance. We have roughly $n^3$ such pairs where each contributes at most $p^3_{\rm ver} \cdot p^6_{\rm e}$. This is divided by $\left(n^2 \cdot p^2_{\rm ver} \cdot p^4_{\rm e}\right)^2$. Hence, the probability to deviate from the expectation by more than a constant factor is at most $1/(n \cdot p_{\rm ver} \cdot p^2_{\rm e}) < 1/(\sqrt{\alpha n})$.
By Chebyshev's inequality, for any non-negative random variable $X$,
\begin{equation}
    \Pr\left[\left|X - E(X)\right| \geq (E(X)/2)\right] \leq \frac{4V(X)}{E^2(X)}
\end{equation}
Taking $X = \sum_{(a, b)\in {V_2 \choose 2}} Z_{a, b}$ and defining $\bar{Z}_{a, b} = Z_{a, b} - \mu$ where $\mu = E_{(a, b)\in {V_2 \choose 2}}[Z_{a, b}]$,  
we obtain that 
\begin{equation}
    \Pr\left[\left|X - E(X)\right| \geq (E(X)/2)\right] \leq \frac{4\sigma_{} E[\bar{Z}_{a, b} \bar{Z}_{a', b'}]}{n^4\mu^2}
\end{equation}

\begin{claim}
Let $v$ be a vertex with $\Gamma_2(v) = s$ and $\ell$ $C_4$s originated from $v$.
If we take $q = $ length-$2$ random walks from $v$, w.h.c.p. we will not see a $C_4$.
\end{claim}

\begin{proof}
Let $S$ denote the set of length-$2$ paths originated from $v$. 
Let $L \subseteq S$ denote the $2\ell$ paths that reside on the $C_4$s originated from $v$.
For the sake of the analysis, we split the process that samples $q$ length-$2$ random walk from $v$ into two steps.
First, it decides for each sample whether it is a walk from $L$ or $S\setminus L$. The former happens with probability $(2\ell)/s$.
After deciding this for each sample, it picks uniformly a random walk from the corresponding set.
Clearly, this is equivalent to selecting uniformly $q$ random walks.

If we pick $q$ random walks, then on expectation, the number of times we pick a random walk from $L$ is $(2\ell q)/s$.
Let $t$ denote this number.
By Chernoff's bound, with high probability, $t$ is concentrated around the expectation. 
Let $E_1$ denote the event that this is indeed the case.

Given $t$, the probability that we don't see a $C_4$ is $(1-1/s)\cdot (1-2/s)\ldots \cdot (1-(t-1)/s) = \Pi_{j=1}^{t-1}\left(1 - j/s\right)$.
Using $(1-1/x)^x \geq 1/e^2$ \commentt{RL: is this a real inequality?} we obtain that this is at least  $\Pi_{j=1}^{t-1} e^{-2j/s} = e^{-t(t-1)/s}$, which is sufficiently large when $t = o(\sqrt{s})$.
The latter happens with high probability when $q = o(s^{3/2}/\ell)$.
\end{proof}

\begin{claim}
    With high probability for every $v\in V$, $s \geq \alpha d$ and $\ell \leq $ 
\end{claim}